\documentclass[acmsmall,screen,nonacm,11pt]{acmart}
\geometry{letterpaper,margin=1in}
\settopmatter{printacmref=false}
\setcopyright{none}
\pagestyle{plain}
\renewcommand{\bibfont}{\normalfont\normalsize}
\newtheorem{theorem}{Theorem}
\newtheorem{proposition}[theorem]{Proposition}
\newtheorem{lemma}[theorem]{Lemma}
\newcommand{\emp}{\varnothing}
\newcommand{\calA}{\mathcal{A}}
\DeclareMathOperator{\lcm}{lcm}
\title{Tight Mortality Bounds for Image-Bounded Nondeterministic Automata}
\author{Zijian Zeng}
\affiliation{\institution{UCSI University}\city{}\country{Malaysia}}
\email{1002266693@ucsiuniversity.edu.my}
\authorsaddresses{Author's contact: Zijian Zeng, UCSI University,
\href{mailto:1002266693@ucsiuniversity.edu.my}{1002266693@ucsiuniversity.edu.my}}
\begin{abstract}
\normalsize
A nondeterministic finite automaton is $k$-image-bounded if every word maps
every individual state to at most $k$ states. A word is mortal if it maps
the entire state set to the empty set. Kiefer and Ryzhikov showed that an
incomplete $n$-state $k$-image-bounded automaton has a mortal word of length
at most $n^{k+1}$ and asked whether this bound is tight.

We resolve the fixed-parameter tightness question: for every fixed integer
$k\geq2$, the maximum shortest mortal-word length is
$\Theta_k(n^{k+1})$. The lower bound holds for every sufficiently large
$n$, even for binary strongly connected automata. The construction combines
a linear-size queue with a packet distributed among $k$ cyclic clocks.
The clocks realize a long least-common-multiple delay using only the sum of
their lengths in states. A global image invariant preserves the bound $k$
under every word, and an exact potential forces a linear number of disjoint
full-delay events. We give explicit constants, a constructive mortality
argument, and a self-contained proof of the matching known upper bound.
\end{abstract}
\keywords{nondeterministic finite automata, image-bounded automata, mortality,
extremal word length}
\hypersetup{pdfauthor={Zijian Zeng}}
\begin{document}
\maketitle
\thispagestyle{plain}
\clearpage

\section{The question and the result}

An automaton can admit a word that eliminates every possible computation,
yet require a long word to do so. Bounding the number of endpoints reachable
from each state gives a natural hierarchy between deterministic automata
and unrestricted nondeterministic automata. The endpoint bound must hold
after every word: a bound on the number of transitions carrying one letter
does not suffice.

Formally, an NFA is a finite labelled relational semiautomaton
$\calA=(Q,\Sigma,\Delta)$, with no designated initial or final states.
For $X\subseteq Q$ and $w\in\Sigma^*$, write $X\cdot w$ for the union of
the endpoints of all $w$-labelled paths starting in $X$; abbreviate
$\{q\}\cdot w$ to $q\cdot w$. Missing transitions are allowed.
The automaton is \emph{$k$-image-bounded} if
\begin{equation}\label{eq:global}
 |q\cdot w|\leq k\qquad(q\in Q,\ w\in\Sigma^*).
\end{equation}
A word $w$ is \emph{mortal} if $Q\cdot w=\emp$. We call the automaton
\emph{incomplete} when a mortal word exists, and denote the minimum length
of such a word by $\mu(\calA)$. This terminology follows
Kiefer and Ryzhikov~\cite{KR2025}; incompleteness here means existence of a
mortal word, rather than merely existence of a missing transition.

Let $M_k(n)$ be the maximum of $\mu(\calA)$ over incomplete $n$-state
$k$-image-bounded NFAs and finite alphabets. This maximum exists: on a fixed
state set there are finitely many letter relations, and duplicate relations
do not affect shortest lengths. Kiefer and Ryzhikov~\cite[Section~9]{KR2025}
give the bound $M_k(n)\leq n^{k+1}$ and ask whether it is tight.
We answer the fixed-$k$ asymptotic interpretation of that question.

\begin{theorem}\label{thm:main}
For each integer $k\geq2$, set
\[
 A_k=1+\frac{k(k+1)}2,\qquad
 N_k=2k+2A_k k!,\qquad
 c_k=\frac{k!}{(2A_k)^{k+1}}.
\]
For every integer $n\geq N_k$, there exists a binary strongly connected
incomplete $n$-state $k$-image-bounded NFA $\calA$ satisfying
\[
 \mu(\calA)\geq c_k n^{k+1}.
\]
Consequently, for $n\geq N_k$,
\[
 c_k n^{k+1}\leq M_k(n)\leq n^{k+1},
\]
and $M_k(n)=\Theta_k(n^{k+1})$ for every fixed $k\geq2$.
\end{theorem}

The constants and threshold may depend on $k$; the theorem determines the
exponent at every fixed level $k\geq2$. The general subset bound $2^n-1$
limits the joint dependence on $k$ and $n$. Strong connectivity means that
the underlying directed graph is strongly connected. Our lower-bound
examples have this property and a binary alphabet. The quadratic order
at $k=1$ is known.

\subsection{Prior mechanisms and the additional ingredient}

Ryzhikov~\cite[Section~4]{R2024} describes a binary deterministic construction
with a cyclic queue and a physical delay path. Our queue action uses this
known mechanism, up to indexing. His discussion also recalls the sharp
arbitrary-alphabet deterministic mortality bound $n(n+1)/2$, attributed
there to Rystsov. The unrestricted NFA exponential lower bounds in the same
paper do not impose~\eqref{eq:global}.

Repeated deterministic delays also occur in Vorel's appended-tail lemma,
reproduced with proof by Ananichev~\cite[Section~2, Lemma~2.1]{A2017}.
Under its almost-permutation and divisibility hypotheses, adding a tail of
$h$ states to an $r$-state automaton increases its reset threshold by $rh$.
That lemma concerns synchronization to a sink; we use it as prior context
for the delay mechanism, not as a theorem about the present NFA model.

The additional ingredient is a reusable nondeterministic delay. A packet
with one state in each of $k$ cycles has period
$L=\lcm(\ell_1,\ldots,\ell_k)$, although it uses only
$\sum_i\ell_i$ clock states. Pairwise coprime lengths can make $L$ of order
$n^k$ at fixed $k$. The essential issue is to force a linear number of these
delays without allowing resets, queue collisions, or simultaneous forks
to create a shortcut. We arrange for every failed packet to return to a
single queue port, separate from the injective action on the remaining
queue. The potential in Lemma~\ref{lem:lower} then measures exactly how much
progress any word can make. Thus the same clocks are reused, rather than
copied for each queue state. This supplies the additional factor $n$ while
preserving the global endpoint promise.

\section{A queue and a packet of clocks}

Fix integers $k\geq2$, $m\geq2$, and $\ell_i\geq2$ for $1\leq i\leq k$.
The disjoint state sets are
\[
 P=\{p_0,\ldots,p_{m-1}\},\qquad
 C_i=\{c_{i,0},\ldots,c_{i,\ell_i-1}\},\qquad
 Q=P\cup\bigcup_{i=1}^k C_i.
\]
Let $S=\{c_{i,0}:1\leq i\leq k\}$ and
$L=\lcm(\ell_1,\ldots,\ell_k)$. The alphabet is $\{a,b\}$.
Letter $a$ rotates each of $P,C_1,\ldots,C_k$ forward by one position:
\[
 p_j\cdot a=\{p_{j+1\bmod m}\},\qquad
 c_{i,j}\cdot a=\{c_{i,j+1\bmod\ell_i}\}.
\]
The other letter is defined by
\begin{align}
 p_0\cdot b&=S,& p_1\cdot b&=\{p_{m-1}\},\label{eq:bqueue}\\
 p_j\cdot b&=\{p_{j-1}\} &&(2\leq j<m),\notag\\
 c_{i,j}\cdot b&=
 \begin{cases}
  \{p_0\},&0\leq j<\ell_i-1,\\
  \emp,&j=\ell_i-1.
 \end{cases}\label{eq:bclock}
\end{align}
Write $f$ for the permutation of $P\setminus\{p_0\}$ induced by $b$.
For $m=2$ this permutation fixes $p_1$. Call $S\cdot a^t$ a
\emph{packet}; its position in every clock is determined by the same $t$.
The state $c_{i,\ell_i-1}$ is the terminal position of clock $i$.

\begin{proposition}\label{prop:threshold}
The automaton just defined is strongly connected and globally
$k$-image-bounded, with least possible image bound exactly $k$.
It is incomplete and satisfies
\begin{equation}\label{eq:threshold}
 1+(m+1)L\ \leq\ \mu(\calA)\ \leq\ 1+(m+1)L+m^2.
\end{equation}
\end{proposition}

We prove the properties separately. The clock lengths need not be coprime
for Proposition~\ref{prop:threshold}, and no condition on $\gcd(m,L)$ is needed.

\begin{lemma}\label{lem:image}
Every image of a singleton under any word is empty, a singleton, or a packet.
The underlying directed graph is strongly connected.
\end{lemma}
\begin{proof}
The family consisting of $\emp$, all singletons, and all packets is closed
under $a$. By~\eqref{eq:bqueue}--\eqref{eq:bclock}, the $b$-image of a
singleton is empty, a singleton, or $S$. For a packet, each component either
disappears under $b$ or reaches the same state $p_0$. Its image is therefore
empty or $\{p_0\}$. This family is closed under both letters and contains
every initial singleton, so induction on the whole word proves the assertion.
Each member has size at most $k$, and $|p_0\cdot b|=k$ proves equality of
the least admissible bound. This induction covers arbitrarily many forks.

Every queue state reaches $p_0$ by a power of $a$. Every clock state reaches
its position zero by a power of $a$ and then reaches $p_0$ by $b$, since
$\ell_i\geq2$. Conversely, from $p_0$ powers of $a$ reach the entire queue;
$b$ followed by suitable powers of $a$ has a path to any prescribed clock
state. These paths prove strong connectivity.
\end{proof}

\section{Charging every unit of progress}

We examine global images immediately after occurrences of $b$, called
\emph{$b$-boundaries}. At such a boundary the image has the form
\[
 T\cup E_\eta,\qquad T\subseteq P,\quad \eta\in\{0,1\},\qquad
 E_0=\emp,\quad E_1=S.
\]
The representation is unique because queue and clock states are disjoint.
For this image define
\[
 V(T,\eta)=|T|+\eta.
\]
This counts queue states and one unit for an entire packet, rather than
counting the individual states in that packet.

\begin{lemma}\label{lem:lower}
Every mortal word has length at least $1+(m+1)L$.
\end{lemma}
\begin{proof}
Letter $a$ is a permutation of $Q$, and
\begin{equation}\label{eq:first}
 Q\cdot b=P\cup S.
\end{equation}
Indeed, $p_0$ supplies $S$, the other queue states supply
$P\setminus\{p_0\}$, and the nonterminal clock states supply $p_0$.
Thus the potential immediately after the first $b$ is $m+1$.

Consider any interval $a^t b$, $t\geq0$, between consecutive occurrences of
$b$. Suppose its starting image is $T\cup E_\eta$.
Put $U=T\cdot a^t$, let $\alpha$ be the indicator of $p_0\in U$, and let
$\varepsilon$ be the indicator that $\eta=1$ and at least one state in
$S\cdot a^t$ is nonterminal. After this interval the queue part is
\[
 T'=f(U\setminus\{p_0\})\cup
 \begin{cases}\{p_0\},&\varepsilon=1,\\\emp,&\varepsilon=0,\end{cases}
\]
and the packet indicator is $\eta'=\alpha$. The displayed union is disjoint,
and $f$ is injective. Since $a$ permutes $P$, we obtain the exact identity
\begin{equation}\label{eq:potential}
 V(T',\eta')=|T|-\alpha+\varepsilon+\alpha
             =|T|+\varepsilon,
 \qquad V(T,\eta)-V(T',\eta')=\eta-\varepsilon.
\end{equation}
In particular, the representation persists after every $b$.

The drop is either zero or one. It equals one precisely when a packet was
present and every component of that packet reaches a terminal clock position.
Since its starting positions are all zero, this condition is
\[
 t\equiv-1\pmod{\ell_i}\quad\hbox{for every }i,
 \qquad\hbox{equivalently}\qquad L\mid t+1.
\]
Every interval that lowers $V$ therefore has length $t+1\geq L$.
The identity includes the case where a killed packet is replaced by a new
packet at the very same $b$: then $\alpha=1$, and the potential still drops
by exactly one. Consecutive $b$'s and empty queue subsets are also covered.

A mortal word contains $b$, because $a$ alone is a permutation. Delete an
initial run of $a$'s, which leaves $Q$ unchanged, and a terminal run of $a$'s,
which cannot turn a nonempty subset into an empty one. The remaining mortal
word starts and ends with $b$. By~\eqref{eq:first}, its first $b$ gives
potential $m+1$; its last gives zero. Equation~\eqref{eq:potential} forces
exactly $m+1$ unit drops. The intervals between consecutive $b$'s are
disjoint, and each of these drops costs at least $L$ letters. Including the
first $b$ proves the lower bound for the shortened word and hence for the
original word.
\end{proof}

\begin{lemma}\label{lem:mortal}
There is a mortal word of length at most $1+(m+1)L+m^2$.
\end{lemma}
\begin{proof}
Begin with $b$, giving potential $m+1$ and a packet. At any $b$-boundary
with a packet, apply $a^{L-1}b$. By~\eqref{eq:potential}, this lowers the
potential by exactly one, whatever the queue subset is.

If the packet is absent but the potential is positive, choose any active
queue state $p_j$. Take $t\in\{0,\ldots,m-1\}$ with $j+t\equiv0\pmod m$,
and apply $a^t b$. This creates a packet and leaves the potential unchanged:
in~\eqref{eq:potential} we have $\eta=\varepsilon=0$ and $\alpha=1$.
This preliminary fork costs at most $m$ letters. Now apply $a^{L-1}b$ to
lower the potential by one.

After exactly $m+1$ such drops the potential is zero, so the global image
is empty. The first drop needs no preliminary fork, and at most $m$ later
forks are needed. The total length is at most $1+(m+1)L+m^2$.
\end{proof}

Lemmas~\ref{lem:image}--\ref{lem:mortal} prove
Proposition~\ref{prop:threshold}, including an explicit mortality procedure.

\section{Obtaining every sufficiently large state count}

\begin{proof}[Proof of the lower bound in Theorem~\ref{thm:main}]
Fix $k\geq2$, write $K=k!$ and $A=A_k$, and take any $n\geq N_k$.
Choose
\begin{equation}\label{eq:parameters}
 M=K\left\lfloor\frac{n-k}{AK}\right\rfloor,\qquad
 \ell_i=1+iM\ (1\leq i\leq k),\qquad
 m=n-k-(A-1)M.
\end{equation}
Since $AM\leq n-k$, we have $m\geq M$. The floor estimate gives
\[
 M\geq\frac{n-k}{A}-K\geq\frac{n}{2A},
\]
where the second inequality follows from $n\geq2k+2AK$.
Also $(n-k)/(AK)\geq2$, so $M\geq2K\geq2$.
Thus $m\geq2$ and $\ell_i\geq2$, as required. Moreover,
\[
 m+\sum_{i=1}^k\ell_i
 =m+k+\frac{k(k+1)}2M=n.
\]
There is no padding by extra states, and strong connectivity is retained.

The lengths $\ell_i$ are pairwise coprime. To see this, suppose $i<j$
and $d$ divides both $1+iM$ and $1+jM$. Then $\gcd(d,M)=1$, while
$d\mid(j-i)M$, so $d\mid j-i$. If $d>1$, a prime divisor $p$ of $d$
satisfies $p\leq j-i\leq k-1$ and therefore divides $K$ and $M$.
This contradicts $\gcd(d,M)=1$. Hence $d=1$.

It follows that
\[
 L=\prod_{i=1}^k(1+iM)\geq K M^k.
\]
Proposition~\ref{prop:threshold} now yields
\[
 \mu(\calA)\geq1+(m+1)L
 \geq K M^{k+1}
 \geq\frac{k!}{(2A_k)^{k+1}}n^{k+1}.
\]
The construction is binary, strongly connected, and globally
$k$-image-bounded by the same proposition, proving the claimed lower bound
for every $n\geq N_k$.
\end{proof}

\section{The matching upper bound}

For completeness we reproduce the elementary upper argument behind the
bound of Kiefer and Ryzhikov~\cite[Section~9]{KR2025}.
Let $\calA$ be any incomplete $n$-state $k$-image-bounded NFA and put
\[
 D=\sum_{j=1}^{\min(k,n)}\binom nj.
\]
Consider any nonempty image $q\cdot u$ of an original state.
A killing continuation exists, since a global mortal word kills every
subset. Along a shortest killing continuation all the intermediate
nonempty subsets have size at most $k$: they are images of $q$ under
longer words, so~\eqref{eq:global} applies.
No nonempty subset can repeat along this continuation. If it did, deleting
the factor between equal subsets would preserve the final empty image and
shorten the continuation. Counting the nonempty subsets before the final
empty image shows that the continuation has length at most $D$.

List the original states $q_1,\ldots,q_n$. Starting with the empty prefix,
successively append a shortest killing continuation for the current image
of $q_i$, skipping it if that image is already empty. A dead row remains
dead under extensions, so after at most $n$ continuations the global image
is empty. The resulting mortal word has length at most $nD$.

Finally, order $Q$. Encode each nonempty subset of size at most $k$ by
listing its elements increasingly and padding the list to length $k$ with
copies of its maximum. This is an injection into $Q^k$, because the set of
entries recovers the subset. Thus $D\leq n^k$, giving
$\mu(\calA)\leq n^{k+1}$. Together with the lower bound, this completes
Theorem~\ref{thm:main}.

\paragraph{Scope and remaining extremal questions.}
The bounds leave room to improve constants and determine exact extremal
thresholds. They concern witness length and do not establish lower bounds on
algorithmic running time. No unambiguity assumption is imposed, and the result
does not resolve mortality questions restricted to unambiguous automata.
The distinction
between endpoints and path multiplicity is visible in the construction:
there are $k$ different $b^2$-paths from $p_0$ to itself, and $k^r$ such
$b^{2r}$-paths. Their endpoints merge, in agreement with
Lemma~\ref{lem:image}.

\paragraph{Computational checks.}
Exact finite checks regenerated the transition relations, exhausted the
reachable global subset graphs and all singleton subset closures in small
instances, and tested the potential identity and the constructive strategy.
A fresh separately implemented verifier covered 61 instances, checked nine
serialized certificates, and tested 277,168 abstract potential transitions,
including simultaneous packet killing and replacement. The cases include
repeated and noncoprime clock lengths. These finite checks test implementation
and boundary cases and complement the all-parameter proof above.

\paragraph{AI methodology.}
OpenAI Codex contributed substantially to all sections: developing the
construction and proofs, comparing primary literature, implementing exact
set and integer-bitmask checks, and drafting the exposition and LaTeX.
Separate Codex contexts checked the argument and finite computations.
Responsibility for any submission and its content remains with the named author.

\bibliographystyle{ACM-Reference-Format}
\bibliography{references}
\end{document}
